\begingroup
\ArithmeticHead{243}{R}{SCHUMACHER AN GAUSS.}
\SourceLeaf{250b}
\pdfbookmark[1]{[5.] Gauss an Schumacher, 17. September 1808}{x1-correspondence-5}
\NumberedItem{5}
\begin{center}\textsc{Gauss} an \textsc{Schumacher}.\end{center}
\HistoricCitation{[Briefwechsel zwischen \textsc{Gauss} und \textsc{Schumacher} I, Altona 1860, S. 2.]}
\par\indent Vor allen Dingen muss ich Sie, mein theuerster Herr Doctor, um Vergebung bitten, dass ich Ihre verehrte Zuschrift vom $2^{\mathrm{ten}}$ May\textsuperscript{[*)]} so lange unbeantwortet gelassen habe. Die Sünde des Aufschiebens wird so leicht zur Gewohnheit, wenn man öfters mit Arbeiten beschäftigt ist, die den grössten Theil unsrer Zeit in Anspruch nehmen.

\par\indent Was die von Ihnen erwähnte Aufgabe betrifft, so muss ich Ihnen aufrichtig gestehen, dass ich bisher es noch nicht habe über mich gewinnen können, sie zum Gegenstande einer besondern ernstlichen Untersuchung zu machen. Ich habe die Unart, \textls[120]{ein lebhaftes} Interesse bei mathematischen Gegenständen nur da zu nehmen, wo ich sinnreiche Ideenverbindungen und durch Eleganz oder Allgemeinheit sich empfehlende Resultate ahnen darf,
\par\vfill\noindent\rule{.25\linewidth}{.4pt}\par\smallskip
\begingroup\footnotesize\noindent [*) So in der Handschrift; richtig wäre \textls[120]{April}.]\par\endgroup
\smallskip\noindent\hfill\small $31^{*}$\normalsize
\clearpage
\ArithmeticHead{244}{L}{ANALYSIS. BRIEFWECHSEL.}
\SourceLeaf{251}
\noindent und wenn ich offenherzig sprechen soll, muss ich erklären, dass mich so etwas aus jenem Problem nicht angesprochen hat. Es kann wol seyn, dass ich irre, und dass wirklich an \textsc{Pedrayes} Problem mehr ist, als man aus seiner Exposition schliessen kann, aber aller Wahrscheinlichkeit nach lässt sich wol nicht viel von jemand erwarten, der sein Problem so verworren vorträgt, dass man den Sinn nur errathen muss. Übrigens weiss ich kaum, ob sie der Aufl[ösung] des H. Prof. \textsc{Pfaff} nicht Unrecht thun, wenn Sie sagen, dass er die Gleichung nur für einen besonderen Fall integrirt habe: wenn \textsc{Pfaff} den Sinn der Aufgabe errathen hat, so ist die wahre Pointe der Aufgabe nicht das Integriren, sondern das Angeben einer Gleichung zwischen $x$ und $u$, bei der jene Diff[erential]-Gleichung sich algebr[aisch] integriren lässt, und wenn es eine andere Aufl[ösung] gibt, als die des Hrn. \textsc{Pfaff}, so scheint es wird dieselbe bloss in der Aufstellung \textls[120]{einer andern} Relation zwischen $x$ und $u$ bestehen als der einfachsten $x=u$, z. B. vielleicht in einer ähnlichen wie $xx+uu=1$\textsuperscript{[*)]}.

\par\indent Vielleicht wäre ich im Besitz von Wahrheiten, die zur Entscheidung dieser Sache dienen könnten. Mir ist bei der Integralrechnung immer das weit weniger interessant gewesen, wo es nur auf Substituiren, Transformiren etc. kurz auf einen gewissen geschickt zu handhabenden Mechanismus ankommt, um Integrale auf algebraische oder Logarithmische oder Kreisfunctionen zu reduciren, als die genauere tiefere Betrachtung solcher Transcendenten Functionen, die sich auf jene nicht zurückführen lassen. Mit Kreisfunctionen und Logarithmischen wissen wir jetzt umzugehen, wie mit dem 1 mal 1, aber die herrliche Goldgrube, die das Innere der höheren Functionen enthält, ist noch fast ganz Terra Incognita. Ich habe darüber ehemals sehr viel gearbeitet und werde dereinst ein eignes grosses Werk darüber geben, wovon ich bereits in meinen \textit{Disquiss. arithm.} p. 593\textsuperscript{[**)]} einen Wink gegeben habe. Man geräth in Erstaunen über den überschwenglichen Reichthum an neuen höchst interessanten Wahrheiten und Relationen, die dergleichen Functionen darbieten (wohin u. a. auch diejenigen gehören, mit denen die Rectification der Ellipse und Hyperbel zusammenhängt). Es könnte wol sein, dass gerade aus diesen
\par\vfill\noindent\rule{.25\linewidth}{.4pt}\par\smallskip
\begingroup\footnotesize
\noindent [*) Die Handschrift hat $xx+yy=1$; in der Aufgabe des \textsc{Pedrayes} heissen die beiden Veränderlichen $x$ und $u$.]\par
\noindent [**) Art. 335, Werke I, S. 412—413.]\par\endgroup
\clearpage
\ArithmeticHead{245}{R}{GAUSS AN SCHUMACHER.}
\SourceLeaf{252}
\noindent Untersuchungen die Beantwortung der \textsc{Pedrayes}-Aufgabe sich entnehmen liesse, vorausgesetzt, dass sie eine Auflösung zulässt, die wirklich einen Werth hat: allein wenn ich auch klarer sähe, dass die ganze Aufgabe zu etwas führen könnte, als dies bis jetzt der Fall ist, würde ich doch jetzt von dieser Untersuchung abstrahiren müssen, da ich mich erst dann in diese weitaussehende Materie wieder hinein werfen werde, wenn ich an die Ausarbeitung jenes grossen Werks werde denken können. Dazu bin ich aber jetzt noch mit zu vielen andern mir nicht minder interessanten Untersuchungen überhäuft.

\par\indent Sollten Sie des Hrn. \textsc{Pedrayes} Auflösung erhalten haben, so würden Sie mich immer sehr durch die Mittheilung verbinden, und es würde mich gewiss innigst freuen, wenn ich finden würde, dass ich mir eine falsche Vorstellung von seiner Aufgabe gemacht habe. Haben Sie sie freilich bis jetzt noch nicht, so werden Sie sie schwerlich jetzt durch den Marquis \textsc{de Romana} erhalten.

\par\indent Von meinem Werke über die Bewegung der Himmelskörper\textsuperscript{[*)]} sind leider erst 17 Bogen, also etwa ${}^3\! /\! {}_5$ des Ganzen fertig. Ich hoffe, dass auf den Novbr. der Druck vollendet seyn wird.

\par\indent Wird die von Ihnen angekündigte Bearbeitung von \textsc{Carnot}, \textit{Géométrie de Position} [Paris, An XI-1803\textsuperscript{**)}] bald erscheinen?

\par\indent Mit ausgezeichneter Hochachtung habe ich die Ehre zu beharren
\begin{flushright}Ew. Wohlgeboren\par\medskip ergebenster Diener\end{flushright}
\noindent Göttingen, d. 17. September 1808.\hfill\textsc{C. F. Gauss.}
\par\vfill\noindent\rule{.25\linewidth}{.4pt}\par\smallskip
\begingroup\footnotesize
\noindent [*) \textit{Theoria motus corporum coelestium}, Hamburgi 1809, Werke VII, 1906, S. 1 ff.; vergl. in Bezug auf den Druck dieses Werkes die Briefe von \textsc{Gauss} an \textsc{Olbers}, \textit{Wilhelm Olbers, Sein Leben und seine Werke} II, 1, 1900, S. 396, 409, 414, 419, 422, 431.]\par
\noindent [**) Siehe: L. N. M. \textsc{Carnot}, \textit{Geometrie der Stellung}, aus dem Französischen übersetzt mit Anmerkungen von H. C. S. \textsc{Schumacher}, I, II, Altona 1808—1810.]\par\endgroup
\clearpage
\endgroup
